%% TRANSCRIBED (PDF 78-87 = printed II-29..II-38)
%% Sec. 3 "Structure of T_m and applications".
%% NOTE: the first three paragraphs of PDF 78 (II-29) -- the tail of
%% Sec. 2.8, ending "...results of the Appendix to Chapter III." -- are
%% transcribed at the end of ch2_s2.tex; this file starts at the Sec. 3
%% heading.
%% NOTE (Weil citation, II-33): the printed text reads "As shown by Weil
%% (cf. [47], p. 4)".  The main bibliography is exactly [1]..[46]
%% (consecutive, no gaps), so [47] is undefined in the original 1968
%% edition as well -- the sentence is an original-book misprint, almost
%% certainly for [43] (A. Weil, "Adeles and algebraic groups"), cited
%% elsewhere for the same purpose (cf. chap. I, and chap. III, 1.1).
%% Per the re-typeset discipline, the book's own numbering inconsistency
%% is KEPT AS PLAIN TEXT (not \cite'd, so it is not mis-linked to the
%% reprint's [47] = Bloch-Kato); flagged in STATE.md for the audit.

\sechead{3}{Structure of $T_m$ and applications}{STRUCTURE OF $T_m$ AND APPLICATIONS}

\subsechead{3.1}{Structure of $X(T_m)$}

If $w$ is a complex place of $\bar{\Q}$, the completion of $\bar{\Q}$ with
respect to $w$ is isomorphic to $\C$; the decomposition group of $w$ is
thus cyclic of order $2$; its non-trivial element will be denoted by
$c_w$ (the ``Frobenius at the infinite place $w$'').  The $c_w$'s are
conjugate in $G = \Gal(\bar{\Q}/\Q)$; let $C_\infty$ denote their
conjugacy class.  (By a theorem of Artin \cite{artin-collected}, p.~257,
the elements of $C_\infty$ are the only non-trivial elements of finite
order in $G$.)

Let $X(T)$ be the character group of the torus $T$, cf.~\S\ref{sec:II.1.1};
we write $X(T)$ additively and put $Y(T) = X(T) \tens_\Z \Q$.  We
decompose $Y$ as a direct sum $Y = Y^0 \oplus Y^- \oplus Y^+$ of
$G$-invariant subspaces, as follows (cf.~Appendix, \S\ref{sec:II.A.2})
\[
Y^0 = Y^G = \{y \in Y \mid gy = y \text{ for all } g \in G\},
\]
\[
Y^- = \{y \in Y \mid cy = -y \text{ for all } c \in C_\infty\},
\]
and $Y^+$ is a $G$-invariant supplement to $Y^0 \oplus Y^-$ in $Y$; one
proves easily that $Y^+$ is unique, cf.~Appendix, \emph{loc.~cit.}

More explicitly, if $\sigma \in \Gamma$ is an embedding of $K$ into
$\bar{\Q}$, let $[\sigma] \in X(T)$ be the corresponding character of
$T$; the $[\sigma]$'s, $\sigma \in \Gamma$, form a basis of $X(T)$ and
$g.[\sigma] = [g.\sigma]$ if $g \in G$.  The space $Y^0$ is generated by
the norm element $\sum\limits_{\sigma \in \Gamma} [\sigma]$, and its
$G$-invariant supplement is
\[
Y^- \oplus Y^+ = \Bigl\{\sum_{\sigma \in \Gamma} b_\sigma[\sigma] \mid
b_\sigma \in \Q,\ \sum_{\sigma \in \Gamma} b_\sigma = 0\Bigr\}.
\]
Hence, any character $\chi \in X(T)$ can be written in the form
\begin{equation}
\tag{$*$}\label{eq:II.3.1.star}
\chi = a \sum_{\sigma \in \Gamma} [\sigma] + \sum_{\sigma \in \Gamma}
b_\sigma[\sigma], \qquad a, b_\sigma \in \Q,\ \sum b_\sigma = 0,\
a + b_\sigma \in \Z.
\end{equation}
(In particular, we see that $da \in \Z$ where $d = [K : \Q]$.)  The
subspace $Y^-$ can now be described as follows
\[
Y^- = \Bigl\{\sum b_\sigma[\sigma] \mid b_\sigma \in \Q,\
\sum b_\sigma = 0,\ b_{c\sigma} = -b_\sigma \text{ for all }
c \in C_\infty \text{ and } \sigma \in \Gamma\Bigr\}.
\]

On the other hand, the projection $T \to T_m$ defines an injection of
$X(T_m)$ into $X(T)$; we identify $X(T_m)$ with its image under this
injection.

\begin{proposition}\label{prop:II.3.1.1}
$X(T_m) \tens_\Z \Q = Y^0 \oplus Y^-$.
\end{proposition}

This follows from Appendix, \S\ref{sec:II.A.2}.

\begin{corollary}[1]\label{cor:II.3.1.1}
The character group $X(T_m)$ is a sublattice of finite index of
$X(T) \cap (Y^0 \oplus Y^-)$.
\end{corollary}

\begin{corollary}[2]\label{cor:II.3.1.2}
If $\chi \in X(T_m)$ is written in the form \eqref{eq:II.3.1.star}, then
$2a \in \Z$.
\end{corollary}

In fact, given $c \in C_\infty$ and $\sigma \in \Gamma$, we have
\[
2a = 2a + b_\sigma + b_{c\sigma} = (a + b_\sigma) + (a + b_{c\sigma})
\in \Z.
\]

\subsechead{3.2}{The morphism $j^*\colon\Gm \to T_m$}

We have seen that any character $\chi \in X(T_m)$ can be written in the
form
\[
\chi = a \sum_{\sigma \in \Gamma} [\sigma] + \sum_{\sigma \in \Gamma}
b_\sigma[\sigma]
\]
with $a, b_\sigma \in \Q$, $\sum b_\sigma = 0$, $2a \in \Z$.  Hence
$\chi \mapsto 2a$ defines a homomorphism
$j\colon X(T_m) \to X(\Gm) = \Z$ and we obtain by duality a morphism of
algebraic groups $j^*\colon \Gm \to T_m$.  If
$\phi_o\colon S_m \to \GL_{V_o}$ is a representation of $S_m$, we obtain
by composition with $j^*$ a morphism of algebraic groups
$\Gm \to \GL_{V_o}$.  This representation of $\Gm$ defines (and is
defined by) a grading $V_o = \sum\limits_{i \in \Z} V_o^{(i)}$ of $V_o$;
recall that $\Gm$ acts on $V_o^{(i)}$ by means of the character
$i \in \Z = X(\Gm)$.

We say that $V_o$ is \emph{homogeneous of degree} $n$ if
$V_o = V_o^{(n)}$.

\runhead{Remark}

For representations coming from the \textit{l}-adic homology
$H_*(\bar{X})$ of a projective smooth variety $X$, the grading defined
above should coincide with the natural one:
$H_*(\bar{X}) = \sum\limits_i H_i(\bar{X})$.

\runhead{Exercise}

\begin{nlist}
\item Let $\mathrm{N}\colon S_m \to \Gm$ be the morphism defined in
Exercise 2 of \S\ref{sec:II.2.5}.  Show that
$\mathrm{N} \circ j\colon \Gm \to S_m \to \Gm$ is
$\lambda \mapsto \lambda^2$.  Show that any morphism $S_m \to \Gm$ is
equal to $\epsilon \mathrm{N}^n$, where $\epsilon$ is a character of
$C_m$ with values in $\{\pm 1\}$ and $n \in \Z$.

\item Let $\phi\colon S_m \to \GL_{V_o}$ be a linear representation of
$S_m$.  Assume $\phi$ is homogeneous of degree $d$, and put
$h = \dim V_o$.
\begin{alist}
\item Show that $dh$ is even (apply Exerc.~1 to
$\det(\phi)\colon S_m \to \Gm$).

\item Prove that there exists on $V_o$ a positive definite quadratic
form $Q$ such that
\[
Q(\rho(x)y) = \mathrm{N}(x)^d Q(y)
\]
for any $y \in V_o$ and any $x \in S_m(\Q)$.  [Let $H$ be the kernel of
$\mathrm{N}\colon S_m \to \Gm$.  Using the fact that $H(\R)$ is compact,
prove the existence of a positive definite quadratic form $Q$ on $V_o$
invariant by $H$; then note that $S_m$ is generated by $H$ and
$j^*(\Gm)$.]
\end{alist}
\end{nlist}

\subsechead{3.3}{Structure of $T_m$}

We need first some notations:

Let $H_c$ be the closed subgroup of $G = \Gal(\bar{\Q}/\Q)$ generated by
$C_\infty$ (cf.~\ref{sec:II.3.1}).  There is a unique continuous homomorphism
$\epsilon\colon H_c \to \{\pm 1\}$ such that $\epsilon(c) = -1$ for all
$c \in C_\infty$.  Indeed the unicity of $\epsilon$ is clear, and one
proves its existence by taking the restriction to $H_c$ of the
homomorphism $G \to \{\pm 1\}$ associated with an imaginary quadratic
extension of $\Q$.  We let $H = \Ker(\epsilon)$.  The groups $H$ and
$H_c$ are closed invariant subgroups of $G$, and $(H : H_c) = 2$.

Let now $K$ be, as before, a finite extension of $\Q$; we identify it
with a subfield of $\bar{\Q}$; let $G_K = \Gal(\bar{\Q}/K)$ be the
corresponding subgroup of $G$.  The field $K$ is \emph{totally real} if
and only if all the elements $c$ of $C_\infty$ act trivially on $K$,
i.e.~if and only if $G_K$ contains $G_c$.  Hence, there exists a
\emph{maximal totally real subfield} $K_o$ of $K$, whose Galois group is
$G_{K_o} = G_K \cdot H_c$.  We let $K_1$ be the field corresponding to
$G_K \cdot H$.  We have
\[
K_o \subset K_1 \subset K \qquad\text{and}\qquad [K_1 : K_o] = 1
\text{ or } 2.
\]

As shown by Weil (cf.~[47], p.~4) the fields $K_o$ and
$K_1$ are closely connected to the groups $T_m$ relative to $K$.
Indeed, if $\chi = \sum b_\sigma[\sigma]$ is an element of the group
denoted by $Y^-$ in \ref{sec:II.3.1}, we have $b_{c\sigma} = -b_\sigma$ for all
$c \in C_\infty$.  If $h = c_1 \ldots c_n$, this gives
\[
b_{h\sigma} = (-1)^n b_\sigma = \epsilon(h) b_\sigma
\]
and by continuity the same holds for all $h \in H_c$.  One deduces from
this:

\begin{proposition}\label{prop:II.3.3.1}
The norm map defines an isomorphism of the space $Y^0_{K_1}$, relative
to $K_1$, onto the space $Y^-_K$, relative to $K$.
\end{proposition}

More precisely, if $\chi_1 = \sum b_{\sigma_1}[\sigma_1]$ belongs to
$Y^-_{K_1}$, where $\sigma_1 \in \Gamma_{K_1}$, the image of $\chi_1$ by
the norm map is
\[
\mathrm{N}^*_{K_1/K_o}(\chi_1) = \sum b_{\sigma/K_1}[\sigma], \qquad
\sigma \in \Gamma_K,
\]
where $\sigma/K_1$ is the restriction of $\sigma$ to $K_1$.  It is clear
that this map is injective.  Conversely, if
$\chi = \sum b_\sigma[\sigma]$ belongs to $Y^-_K$, we saw above that
$b_{h\sigma} = \epsilon(h)b_\sigma$ for all $h \in H_c$, hence
$b_{h\sigma} = b_\sigma$ for $h \in H$ and of course also for
$h \in H.G_K$.  This shows that $b_\sigma$ depends only on the
restriction of $\sigma$ to $K_1$, and hence that $\chi$ belongs to the
image of the norm map.

\begin{corollary}\label{cor:II.3.3.1}
The tori $T_m$ attached to $K$ and $K_1$ are isogenous to each other.
\end{corollary}

There remains to describe the tori $T_m$ attached to $K_1$.  There are
two cases:

\begin{nlist}
\item $K_1 = K_o$.  In this case, we have $Y^- = 0$ and $T_m$ is
one-dimensional, and isomorphic to $\Gm$.

Indeed, if $\chi = \sum b_\sigma[\sigma]$ belongs to $Y^-$, and
$c \in C_\infty$, we have $b_{c\sigma} = -b_\sigma$ (cf.~\ref{sec:II.3.1}) but also
$b_{c\sigma} = b_\sigma$ since $c \in G_K.H_c = G_K.H$.  This shows that
$b_\sigma = 0$ for all $\sigma$, hence $Y^- = 0$.

\item $[K_1 : K_o] = 2$.  The field $K_1$ is then a totally imaginary
quadratic extension of $K_o$ (and it is the only one contained in $K$,
as one checks readily).  In this case $Y^-$ is of dimension
$d = [K_o : \Q]$ and $T_m$ is $(d+1)$-dimensional.
\end{nlist}

More precisely, the space $Y$ attached to $K_1$ is $2d$-dimensional and
the involution $\sigma$ of $K_1$ corresponding to $K_o$ decomposes $Y$
in two eigenspaces of dimension $d$ each; the space $Y^-$ is the one
corresponding to the eigenvalue $-1$ of $\sigma$.  This is proved by the
same argument as above, once one remarks that all $c \in C_\infty$
induce $\sigma$ on $K_1$.

\runhead{Remark}

In this last case (which is the most interesting one), the torus $T_m$
is isogenous to the product of $\Gm$ by the $d$-dimensional torus kernel
of the norm map from $K_1$ to $K_o$.

\subsechead{3.4}{How to compute Frobeniuses}

Let $\phi$ be a linear representation of $S_m$ of degree $n$.  By
extending the groundfield, the restriction of $\phi$ to $T_m$ can be put
in diagonal form; let $\chi_1, \ldots, \chi_n$ be the $n$ characters of
$T_m$ so obtained and write (in additive notation)
\[
\chi_i = \sum_{\sigma \in \Gamma} n_\sigma(i)[\sigma] \qquad
(n_\sigma(i) \in \Z).
\]
We say that $\chi_i$ is \emph{positive} if all the $n_\sigma(i)$'s are
$\geq 0$.  Let $v \notin \Supp(m)$, and let $F_v \in S_m(\Q)$ be the
corresponding Frobenius element, cf.~\ref{sec:II.2.3}.  Since $C_m = S_m/T_m$ is
finite, there exists an integer $N \geq 1$ such that
$F_v^N \in T_m(\Q)$.  If $\mathfrak{p}_v$ is the prime ideal of $v$,
this means that there exists $\alpha \in K^*$, with
$\mathfrak{p}_v^N = (\alpha)$, $\alpha \equiv 1 \bmod m$, and $\alpha > 0$
at all real places of $K$.

\begin{proposition}[1]\label{prop:II.3.4.1}
The eigenvalues of $\phi(F_v^N)$ are the numbers
\[
\chi_i(\alpha) = \prod_\sigma \sigma(\alpha)^{n_\sigma(i)} \qquad
(i = 1, \ldots, n).
\]
\end{proposition}

This is trivial by construction, because $F_v^N$ is the image of
$\alpha$ under $T(\Q) \to T_m(\Q)$.

\begin{corollary}[1]\label{cor:II.3.4.1}
The eigenvalues of $\phi(F_v)$ are $\{p_v\}$-units (i.e.~they are units
at all places of $\bar{\Q}$ not dividing $p_v$).
\end{corollary}

\begin{corollary}[2]\label{cor:II.3.4.2}
Let $z_1, \ldots, z_n$ be the eigenvalues of $\phi(F_v)$, indexed so that
$z_i^N = \chi_i(\alpha)$.  Let $w$ be a place of $\bar{\Q}$ dividing
$p_v$, normalized so that $w(p_v) = v(p_v) = e_v$.  Then
\[
w(z_i) = \sum_{\substack{\sigma \in \Gamma \\ w \circ \sigma = v}}
n_\sigma(i).
\]
\end{corollary}

We have
$w(z_i^N) = w\bigl(\prod\limits_{\sigma \in \Gamma}
\sigma(\alpha)^{n_\sigma(i)}\bigr) = \sum\limits_{\sigma \in \Gamma}
n_\sigma(i)\, w \circ \sigma(\alpha)$, and
\[
w \circ \sigma(\alpha) = 0 \quad\text{if}\quad w \circ \sigma \neq v
\]
\[
w \circ \sigma(\alpha) = N \quad\text{if}\quad w \circ \sigma = v,
\]
since $(\alpha) = \mathfrak{p}_v^N$.  Hence the result.

\begin{corollary}[3]\label{cor:II.3.4.3}
Let $\ell$ be a prime number and let
$\phi_\ell\colon \Gal(\bar{K}/K) \to \Aut(V_\ell)$ be the $\ell$-adic
representation of $K$ associated to $\phi$.  Then $\phi_\ell$ is
integral (cf.~chap.~\ref{chap:I}, \S\ref{sec:I.2.2}) if and only if all
the characters $\chi_i$ occurring in $\phi$ are positive.
\end{corollary}

\begin{proof}[Proof of Corollary~\ref{cor:II.3.4.3}]
Assume first the $\chi_i$'s are positive.  Let $v \notin \Supp(m)$ and
let $z_1, \ldots, z_n$ be the corresponding eigenvalues of $\phi(F_v)$
as in Corollary~\ref{cor:II.3.4.2}.  Corollaries~\ref{cor:II.3.4.1} and~\ref{cor:II.3.4.2} show that the $w(z_i)$ are
positive for all valuations $w$ of $\bar{\Q}$; hence the $z_i$ are
integral over $\Z$.  Hence the $\phi_\ell$'s are integral.

Conversely, assume $\phi_\ell$ is integral for some $\ell$.  There exists
a finite subset $S'$ of $\Sigma_K$, containing $\Supp(m)$, such that if
$v \notin S'$, the eigenvalues of $\phi(F_v)$ are integral.  Choose a
prime number $p$ which splits completely in $K$ and is such that
$p_v = p$ implies $v \notin S'$.  Let $w$ be a valuation of $\bar{\Q}$
dividing $p$.  The valuations $w \circ \sigma$, $\sigma \in \Gamma$, are
pairwise inequivalent.  Let $\sigma \in \Gamma$; and let $v$ be the
normalized valuation of $K$ equivalent to $w \circ \sigma$ so that
$\lambda v = w \circ \sigma$ for some $\lambda > 0$.  Let
$z_1, \ldots, z_n$ be the eigenvalues of $\phi(F_v)$.  By Corollary~\ref{cor:II.3.4.2},
$w(z_i) = \lambda\, n_\sigma(i)$.  Since the $z_i$ are integral, this
shows that the $n_\sigma(i)$'s are all positive.
\end{proof}

\begin{proposition}[2]\label{prop:II.3.4.2}
Let $v \notin \Supp(m)$ and let $\chi$ be a character of $S_m$.  Let
$\chi_T \in X(T_m)$ be the restriction of $\chi$ to $T_m$ and let
$i = j(\chi_T)$ be the integer defined in \ref{sec:II.3.2}.  Then, for any
archimedian absolute value $\omega$ of $\bar{\Q}$ extending the usual
absolute value of $\Q$, we have
\[
\omega(\chi(F_v)) = (\mathrm{N}v)^{i/2}.
\]
\end{proposition}

\begin{proof}
If $\chi = a \sum\limits_{\sigma \in \Gamma} [\sigma] +
\sum\limits_{\sigma \in \Gamma} b_\sigma[\sigma]$ as in \ref{sec:II.3.1}, we have
\[
\omega(\chi(F_v))^N = \omega(\chi(F_v^N)) = \prod_\sigma
\omega \circ \sigma(\alpha)^a \cdot \prod_\sigma
\omega \circ \sigma(\alpha)^{b_\sigma},
\]
and
$\prod\limits_\sigma \omega \circ \sigma(\alpha)^a
= \omega(\mathrm{N}(\alpha))^a = (\mathrm{N}v)^{aN}
= (\mathrm{N}v)^{iN/2}$, where $i = 2a$.  It remains to show that
$x = \prod\limits_\sigma \omega \circ \sigma(\alpha)^{b_\sigma}$ is equal
to $1$.  Let $c = c_\omega$ be the ``Frobenius'' attached to $\omega$
(cf.~\ref{sec:II.3.1}).  Since $b_\sigma + b_{c\sigma} = 0$, we have $x \cdot y = 1$
with $y = \prod\limits_\sigma \omega \circ \sigma(\alpha)^{b_{c\sigma}}$.
But $y = \prod\limits_\tau \omega \circ c \circ \tau(\alpha)^{b_\tau}$,
and, since $\omega \circ c = \omega$, we have $y = x$, hence $x^2 = 1$,
and $x = 1$, since $x > 0$.
\end{proof}

\runhead{Exercises}

\begin{nlist}
\item Check the product formula for the eigenvalues of the
$\phi(F_v)$.  (Use Cor.~1 and 2 to Prop.~1 and Prop.~2.)

\item Show that Prop.~2 and Cor.~1 and 2 to Prop.~1 determine the
eigenvalues of the $\phi(F_v)$'s up to multiplication by roots of unity.

\item (Generalization of Cor.~1 to Prop.~1).  Let $(\rho_\ell)$ be a
strictly compatible system of rational $\ell$-adic representations, with
exceptional set $S$ (cf.~chap.~\ref{chap:I}, \S\ref{sec:I.2.3}).  Show
that, for any $v \in \Sigma_K - S$, the eigenvalues of
$F_{v,\rho_\ell}$, $\ell \neq p_v$, are $p_v$-units.
\end{nlist}
